\documentclass[11pt]{article}
\usepackage{acl}

\usepackage{times}
\usepackage{latexsym}
\usepackage[T1]{fontenc}
\usepackage[utf8]{inputenc}
\usepackage{microtype}
\usepackage{inconsolata}
\usepackage{url}
\usepackage{booktabs}
\usepackage{amsmath}
\usepackage{amssymb}
\usepackage{enumitem}
\Urlmuskip=0mu plus 1mu
\emergencystretch=3em

\title{Predicting Market Outcomes from Earnings Call Transcripts:\\Finetuning, Retrieval, and Prompting}

\author{
Surya Suresh \\
\And
Dheeraj Gosula \\
\And
Nikhil Kumar \\
}

\begin{document}
\maketitle

\begin{abstract}
We finetune an LLM to predict stock direction and earnings surprise directly from earnings call transcripts, then test whether retrieval-augmented generation and structured prompting (few-shot examples with chain-of-thought reasoning) further improve results. We train separate models for each prediction task and evaluate all combinations of the three factors across eight conditions per task. We also compare predictions from the Q\&A section versus prepared remarks and test robustness across market periods. Preliminary results with a FinBERT sentiment baseline show near-chance prediction, establishing the floor that our approach needs to beat.
\end{abstract}

\section{Introduction}

Earnings call transcripts are a primary source of information for investors. Companies report financial results and forward guidance, and analysts press management with questions. This language contains clues about future stock movement and earnings performance, but making accurate predictions from it is difficult. Prior work has used sentiment analysis and domain-pretrained models with results generally only slightly above chance \citep{Qin2019}.

We take a direct approach: we finetune an LLM on past transcript-outcome pairs to predict stock direction (up or down) and earnings surprise (beat or miss). We then test whether two additional techniques improve results. Retrieval-augmented generation (RAG) gives the model context from prior transcripts at inference time, such as the same company's previous quarter, so it can detect changes in tone or guidance. Structured prompting gives the model few-shot examples of past transcripts with known outcomes and asks it to reason step-by-step before predicting. We train separate models for each prediction task and test every combination of the three factors, giving us eight conditions per task. We also test whether the Q\&A section is more predictive than the prepared remarks, and whether results hold across different market periods.

\section{Related Work}

Financial sentiment analysis has progressed from the Loughran-McDonald dictionary \citep{Loughran2011} to domain-pretrained models like FinBERT \citep{Araci2019}. Earnings call language has been shown to predict volatility \citep{Qin2019}, but transcripts are usually treated as flat text without separating prepared remarks from Q\&A. RAG has been studied in open-domain settings \citep{Lewis2020} but not applied to financial prediction from transcripts. Chain-of-thought prompting \citep{Wei2022} has improved LLM reasoning on many tasks but has not been tested for market prediction. We combine finetuning, retrieval, and prompting in a single experiment to measure what helps and what does not.

\section{Experiment Design}

\subsection{Overview}

For each prediction task (stock direction and earnings surprise), we run eight conditions formed by toggling three factors on or off: finetuning, RAG, and structured prompting. This ranges from a bare zero-shot baseline (all off) to a fully augmented model (all on). We train a separate finetuned model for each task, since the patterns that predict stock direction and earnings surprise are likely different.

\subsection{Finetuning}

We finetune the base LLM (7B parameters) using LoRA \citep{Hu2022} on past transcript-outcome pairs. For stock direction, the training examples are (transcript, ``up'') or (transcript, ``down''), based on whether the stock closed higher or lower over the three trading days after the call. For earnings surprise, the examples are (transcript, ``beat'') or (transcript, ``miss''), based on whether reported EPS exceeded or fell short of analyst consensus. The full dataset contains tens of thousands of transcripts with computable labels. Our preliminary pipeline has processed 394 so far, and we will expand this to thousands for the final experiments. We use LoRA to keep the number of trainable parameters small and reduce overfitting risk.

\subsection{Retrieval-Augmented Generation}

At inference time, the model receives passages retrieved from our transcript corpus. For a given earnings call, we retrieve passages from the same company's prior quarter transcript so the model can compare how tone, guidance, and analyst concerns have changed. Retrieval uses BM25 or dense retrieval with sentence-transformer embeddings \citep{Reimers2019} indexed in FAISS. The retrieved passages are prepended to the current transcript in the prompt.

\subsection{Structured Prompting}

The model receives few-shot examples of past transcripts with their known outcomes (e.g., two where the stock went up and two where it went down) and is asked to reason step-by-step before predicting. The few-shot examples teach the model what patterns to look for, and the chain-of-thought step forces it to articulate its reasoning before committing to an answer.

\subsection{Conditions}

For each prediction target, we run eight conditions: the base model alone (zero-shot), each factor individually (finetuned only, RAG only, prompting only), each pair (finetuned + RAG, finetuned + prompting, RAG + prompting), and all three combined. The finetuned conditions use a task-specific model and are run independently for each target. The non-finetuned conditions use the same base model and predict both targets.

\subsection{Evaluation}

In every condition, the model receives a transcript and predicts stock direction (up/down) or earnings surprise (beat/miss). For example, the prompt might read: ``Given the following earnings call transcript, predict whether the stock price will go up or down in the next three trading days. Answer only `up' or `down'.'' In conditions with structured prompting, the prompt additionally includes few-shot examples and asks for step-by-step reasoning. We measure prediction quality with accuracy, balanced accuracy, and AUC-ROC, with bootstrap confidence intervals and permutation tests for significance.

\subsection{Section Comparison}

We test whether the Q\&A section is more predictive than the prepared remarks by running prediction with only the Q\&A versus only the prepared remarks as input. The unscripted Q\&A may contain more candid language than the scripted remarks \citep{Qin2019}.

\subsection{Temporal Robustness}

We evaluate on three eras: pre-COVID (2015--2019), COVID (2020--2021), and post-COVID (2022--2024), using only preceding transcripts for finetuning and few-shot examples. This tests whether results hold across different market conditions.

\section{Data}

We use the Bose345/sp500\_earnings\_transcripts dataset, which contains tens of thousands of earnings calls from 2010--2024. Stock direction labels come from \texttt{yfinance} (up/down based on closing price movement over three post-call trading days) and earnings surprise labels from \texttt{yahooquery} (beat/miss based on reported EPS vs.\ analyst consensus). All splits are chronological to prevent data leakage. Transcripts are split into prepared remarks and Q\&A sections using rule-based heuristics.

We compare against three baselines that do not use an LLM and use a different prediction method. Majority class always predicts the most common label. The Loughran-McDonald baseline \citep{Loughran2011} counts positive and negative financial words in the transcript and feeds that score into a logistic regression classifier to predict direction. FinBERT \citep{Araci2019}, a sentiment classifier pretrained on financial text, similarly feeds its sentiment scores into logistic regression. These baselines use a simpler pipeline than our LLM conditions (sentiment score $\rightarrow$ classifier, rather than direct prediction from the transcript), but they are useful reference points: if an LLM cannot outperform simple word counting, then the added complexity is not justified.

\section{Work Plan}

In the first phase, we build the data pipeline: clean transcripts, segment into prepared remarks and Q\&A, compute labels, and create chronological splits. In the second phase, we finetune two models (one per task) using LoRA, build the retrieval index, and design the prompts. In the third phase, we run all conditions for both tasks, run the section comparison and temporal robustness checks, and write up findings.

\section{Progress Report}

\subsection{What We Have Built So Far}

We have completed the data pipeline and initial infrastructure. The code is a Python package under \texttt{src/}, running on the OSC Pitzer cluster with V100 GPUs via Slurm (PyTorch 2.3.1, CUDA 11.8, Transformers 4.44.2). The pipeline loads transcripts, splits them into prepared remarks and Q\&A using regex patterns, and computes stock direction labels from closing prices via \texttt{yfinance}. As an initial test, we processed 500 recent transcripts, of which 493 yielded valid labels after excluding delisted tickers (394 train, 99 test). For the final experiments, we will expand to thousands of transcripts from the full corpus, which will also reduce the current label imbalance (309 up vs.\ 184 down) by covering more varied market periods.

\subsection{Preliminary Results}

As an initial baseline, we tested FinBERT \citep{Araci2019} on stock direction prediction. It scores each transcript's sentiment as $P(\text{positive}) - P(\text{negative})$, chunking into 150-word windows, and we fed those scores into a logistic regression classifier. We tested two variants: one score over the full transcript, and separate scores for prepared remarks and Q\&A.

\begin{table}[t]
\centering
\small
\begin{tabular}{@{}lc@{}}
\toprule
\textbf{Condition} & \textbf{AUC-ROC} \\
\midrule
Majority class & 0.500 \\
FinBERT zero-shot & 0.488 \\
FinBERT zero-shot (section-split) & 0.512 \\
\bottomrule
\end{tabular}
\caption{Baseline stock direction AUC on 99 held-out examples. Earnings surprise prediction is not yet implemented.}
\label{tab:preliminary}
\end{table}

Both produce near-chance AUC (0.49 and 0.51), which is expected since FinBERT produces only a generic sentiment score. Management prepared remarks tend to be positive regardless of actual performance \citep{Qin2019}, so aggregate sentiment is uninformative. The section-split variant shows a small gain, consistent with the Q\&A containing more candid language. These baselines establish the floor that our finetuned LLM needs to beat.

\subsection{What Remains}

In the first week, we finetune two models (one for stock direction, one for earnings surprise) using LoRA, build the earnings surprise labels, and design the prompts. In the second week, we build the retrieval index and run all conditions for both tasks. In the third week, we run the section comparison and temporal robustness checks, analyze results, and write the final report.

\section*{References}
\begingroup
\renewcommand{\section}[2]{}
\begin{thebibliography}{}

\bibitem[Araci(2019)]{Araci2019}
Dogu Araci. 2019.
FinBERT: Financial Sentiment Analysis with Pre-trained Language Models.
\textit{arXiv preprint arXiv:1908.10063}.

\bibitem[Hu et al.(2022)]{Hu2022}
Edward Hu et al. 2022.
LoRA: Low-Rank Adaptation of Large Language Models.
In \textit{ICLR}.

\bibitem[Lewis et al.(2020)]{Lewis2020}
Patrick Lewis et al. 2020.
Retrieval-Augmented Generation for Knowledge-Intensive NLP Tasks.
In \textit{NeurIPS}.

\bibitem[Loughran and McDonald(2011)]{Loughran2011}
Tim Loughran and Bill McDonald. 2011.
When Is a Liability Not a Liability?
\textit{Journal of Finance}, 66(1):35--65.

\bibitem[Qin and Yang(2019)]{Qin2019}
Yao Qin and Yi Yang. 2019.
What You Say and How You Say It Matters: Predicting Stock Volatility Using Verbal and Vocal Cues.
In \textit{ACL}, pages 390--401.

\bibitem[Reimers and Gurevych(2019)]{Reimers2019}
Nils Reimers and Iryna Gurevych. 2019.
Sentence-BERT: Sentence Embeddings using Siamese BERT-Networks.
In \textit{EMNLP}, pages 3982--3992.

\bibitem[Wei et al.(2022)]{Wei2022}
Jason Wei et al. 2022.
Chain-of-Thought Prompting Elicits Reasoning in Large Language Models.
In \textit{NeurIPS}.

\end{thebibliography}
\endgroup

\end{document}
